\documentclass[11pt]{article}
\usepackage[utf8]{inputenc}
\usepackage{amsmath}
\usepackage{amsfonts}
\usepackage{booktabs}
\usepackage{geometry}
\usepackage{listings}
\usepackage{xcolor}
\usepackage{hyperref}

\geometry{margin=1in}

\title{The War with Mars: From Kepler's Geometry to AI and Bayesian Inference}
\author{Historical and Computational Analysis}
\date{\today}

\definecolor{codegreen}{rgb}{0,0.6,0}
\definecolor{codegray}{rgb}{0.5,0.5,0.5}
\definecolor{codepurple}{rgb}{0.58,0,0.82}
\definecolor{backcolour}{rgb}{0.95,0.95,0.92}

\lstset{
    backgroundcolor=\color{backcolour},   
    commentstyle=\color{codegreen},
    keywordstyle=\color{blue},
    numberstyle=\tiny\color{codegray},
    stringstyle=\color{codepurple},
    basicstyle=\small\ttfamily,
    breakatwhitespace=false,         
    breaklines=true,                 
    captionpos=b,                    
    keepspaces=true,                 
    numbers=left,                    
    numbersep=5pt,                  
    showspaces=false,                
    showstringspaces=false,
    showtabs=false,                  
    tabsize=2
}

\begin{document}

\maketitle

\section{Historical Calibration: Kepler and Brahe}
Johannes Kepler inherited the observational data of Tycho Brahe, which achieved a then-unprecedented precision of approximately $2'$ (arcminutes) \cite{1.2.7}. Kepler's task was to calibrate a model to fit the heliocentric coordinates of Mars. 

Initially, Kepler utilized the \textbf{Vicarious Hypothesis}---an eccentric circle with an equant point. While this model predicted longitudes within $2'$, it failed to predict distances correctly \cite{1.2.8}. The ultimate "calibration" was driven by a systematic \textbf{8-arcminute error} in the circular model, which led Kepler to abandon perfect circles in favor of the ellipse \cite{1.2.1, 1.2.8}.

\section{The Parameterized Mathematical Model}
Kepler's model effectively optimized a polar representation of planetary motion. In modern notation, the parameterized model he used is:

\begin{equation}
r(\theta) = \frac{a(1 - e^2)}{1 + e \cos(\theta)}
\end{equation}

Where the primary parameters for calibration include:
\begin{itemize}
    \item \textbf{$a$ (Semi-major axis):} The scale of the orbit (approx. 1.524 AU for Mars) \cite{1.2.1}.
    \item \textbf{$e$ (Eccentricity):} The degree of elongation (approx. 0.0934 for Mars) \cite{1.5.7}.
    \item \textbf{$\theta$ (True Anomaly):} The angular position relative to perihelion.
    \item \textbf{$M$ (Mean Anomaly):} The time-varying parameter related to $E$ (Eccentric Anomaly) via $M = E - e \sin E$ \cite{1.2.1}.
\end{itemize}

\section{Modern Statistical Recovery}
Modern work treats Kepler's discovery as a problem of inference and automated discovery.

\subsection{Bayesian and Frequentist Approaches}
In a \textbf{Frequentist} framework, researchers apply non-linear least squares to the Rudolphine Tables, minimizing residuals between observed and predicted positions \cite{1.5.7}. In a \textbf{Bayesian} framework, Markov Chain Monte Carlo (MCMC) sampling is used to determine the posterior distribution of $e$ and $a$. 

Below is a Python snippet illustrating a Bayesian recovery using the \texttt{emcee} library:

\begin{lstlisting}[language=Python]
import numpy as np
import emcee

# Kepler's Elliptical Model
def model(theta, e, a=1.524):
    return (a * (1 - e**2)) / (1 + e * np.cos(theta))

# Bayesian Probability Functions
def log_likelihood(e, theta, r, sigma):
    if not (0.0 < e < 1.0): return -np.inf
    r_pred = model(theta, e)
    return -0.5 * np.sum(((r - r_pred) / sigma)**2)

# MCMC Sampler Setup
nwalkers, ndim = 32, 1
p0 = np.random.rand(nwalkers, ndim)
sampler = emcee.EnsembleSampler(nwalkers, ndim, log_likelihood, args=(theta_obs, r_obs, 0.008))
sampler.run_mcmc(p0, 1000)
\end{lstlisting}

\section{AI and Symbolic Regression}
Modern AI techniques, such as \textbf{Symbolic Regression (SR)}, can rediscover Kepler’s laws from raw data without prior geometric assumptions \cite{1.3.1, 1.5.4}.

\subsection{AI Feynman and Pareto Frontiers}
Tools like \textit{AI Feynman} (MIT) search for the most "parsimonious" equation by exploring a \textbf{Pareto Frontier}---balancing model accuracy against mathematical complexity \cite{1.5.9}. The AI identifies the ellipse because it provides a massive reduction in complexity compared to epicyclic models while maintaining high accuracy \cite{1.5.8}.

\subsection{Noise Handling: Neural Smoothing}
To handle the $2'$ jitter in Brahe's data, modern AI uses \textbf{Neural Smoothing} \cite{1.3.5}. A neural network is first trained to interpolate the noisy data; the SR algorithm then queries this smoothed "denoised" manifold to extract the symbolic law $r(\theta)$ \cite{1.3.5}.

\section{Comparison Summary}
\begin{center}
\begin{tabular}{lll}
\toprule
\textbf{Feature} & \textbf{Kepler (1605)} & \textbf{Modern AI/Stats} \\ \midrule
Optimization & Manual Iteration & MCMC / Gradient Descent \\
Error Metric & Maximum Deviation & Log-Likelihood / MDL \cite{1.3.5} \\
Model Search & Geometric Intuition & Symbolic Regression \cite{1.5.8} \\
Noise Filtering & 8-minute threshold & Neural Interpolation \cite{1.3.5} \\ \bottomrule
\end{tabular}
\end{center}

\begin{thebibliography}{9}
\bibitem{1.2.1} Galileo Project, "Johannes Kepler and Tycho Brahe," University of Virginia.
\bibitem{1.2.7} "Kepler's Discovery of the Elliptical Orbit of Mars," St. John's Digital Archives.
\bibitem{1.3.5} Qeios, "Kepler: The Pioneer of Data Science and AI," 2025.
\bibitem{1.5.7} arXiv:2312.09766, "Celestial Machine Learning," 2023.
\bibitem{1.5.8} Wikipedia, "Symbolic Regression."
\bibitem{1.5.9} arXiv:2312.09766, "Emulating Kepler's discovery with AI Feynman."
\end{thebibliography}

\end{document}
